\documentclass[11pt,a4paper]{article}
\usepackage[utf8]{inputenc}
\usepackage[T1]{fontenc}
\usepackage[english]{babel}
\usepackage{geometry}
\usepackage{amsmath}
\usepackage{hyperref}
\usepackage{parskip}
\usepackage{titlesec}
\usepackage{enumitem}
\usepackage{booktabs}
\usepackage{xcolor}

\geometry{top=2.0cm, bottom=2.0cm, left=2.2cm, right=2.2cm}

\titleformat{\section}{\normalfont\large\bfseries}{}{0em}{}[\titlerule]
\titlespacing{\section}{0pt}{10pt}{5pt}
\titleformat{\subsection}{\normalfont\normalsize\bfseries}{}{0em}{}
\titlespacing{\subsection}{0pt}{6pt}{2pt}

\hypersetup{colorlinks=true, linkcolor=blue!70!black, urlcolor=blue!70!black}

\begin{document}

\begin{center}
  {\LARGE\bfseries Project Description}\\[4pt]
  {\large GN2: Transformer-Based Jet Flavour Tagging in PyTorch}\\[2pt]
  {\small\color{gray}
    Based on: ATLAS Collaboration, \textit{Nature Communications} \textbf{17}, 541 (2026)}
\end{center}
\vspace{4pt}\hrule\vspace{10pt}

% ─── Motivation ──────────────────────────────────────────────────────────────
\section{Motivation}

Hadronic jets are the most abundant physics objects in proton--proton collision
events at the LHC.
Identifying the flavour of a jet --- whether it originates from a $b$-quark,
a $c$-quark, a hadronically-decaying $\tau$-lepton, or a light quark/gluon ---
is crucial for a wide range of ATLAS physics analyses, including measurements
of Higgs boson couplings ($H\to b\bar{b}$, $H\to c\bar{c}$), searches for
new physics beyond the Standard Model, and top-quark studies.
The long lifetime of $b$-hadrons ($\tau\approx 1.5$~ps) causes them to
travel several millimetres before decaying, producing a secondary vertex
significantly displaced from the primary interaction point that can be
reconstructed from inner-tracker information.

Traditional two-stage flavour-tagging algorithms reconstruct displaced
vertices with specialised low-level tools and then combine their outputs in a
high-level multivariate classifier (e.g.\ the DL1d deep neural network).
GN2 replaces this pipeline with a single end-to-end transformer that directly
processes the raw charged-particle tracks associated with the jet, achieving
improvements of a factor 1.5--4 over DL1d across its major applications.

% ─── Dataset & Preprocessing ─────────────────────────────────────────────────
\section{Dataset and Preprocessing (\texttt{preprocess.py}, \texttt{dataset.py})}

The training data consists of ATLAS Monte Carlo simulated $t\bar{t}$ events
(jets with $20 < p_T < 250$~GeV) and hypothetical $Z'$ boson events
(jets with $250 < p_T < 6000$~GeV), stored in HDF5 format.
Each jet provides two level of information:
\textit{jet-level} features ($p_T$, $\eta$) and
\textit{track-level} features for up to 40 associated charged-particle
tracks (transverse and longitudinal impact parameters $d_0$, $z_0$ and their
significances, angular separation variables, and quality flags).
A boolean \texttt{valid} field marks real tracks; padded positions are
zeroed and excluded from attention via a boolean mask.

The preprocessing step (\texttt{preprocess.py}) applies kinematic selection
($20~\text{GeV} < p_T < 250~\text{GeV}$, $|\eta| < 2.5$), performs a
stratified 70/15/15 train/val/test split using \texttt{scikit-learn},
and computes per-feature Z-score normalisation statistics
(with a log-transform applied to $p_T$ before normalisation)
\emph{exclusively on the training set}, avoiding any data leakage into
validation and test sets.
Statistics are saved to \texttt{norm\_stats.json} and reused at inference time.

The custom \texttt{GN2Dataset} and \texttt{\_BatchCollator} classes sort
jet indices before each HDF5 read to guarantee contiguous disk access,
reducing I/O overhead by orders of magnitude compared to per-sample access.

% ─── Architecture ────────────────────────────────────────────────────────────
\section{Model Architecture (\texttt{model.py})}

The GN2 model processes each jet in four sequential stages.

\begin{enumerate}[leftmargin=*, itemsep=3pt, topsep=3pt]
  \item \textbf{Per-track initialiser.}
    Jet features are broadcast to each track slot and concatenated with the
    track features, producing a combined vector of dimension
    $n_\text{jet} + n_\text{track}$ per track.
    A two-layer MLP with hidden size 256 and ReLU activation maps this to
    an initial track representation of dimension 256.

  \item \textbf{Transformer encoder.}
    Four encoder layers, each with 8 attention heads, embedding dimension 256,
    and feed-forward dimension 512, apply multi-head self-attention with
    pre-LayerNorm.
    Padded track slots are excluded from attention via a key-padding mask.

  \item \textbf{Attention pooling.}
    A learned linear query scores each track embedding;
    softmax weights aggregate the embeddings into a single global jet
    representation of dimension 128.
    Fully-masked jets receive uniform weights as a safe fallback.

  \item \textbf{Classification head.}
    A three-hidden-layer MLP ($128 \to 64 \to 32 \to 4$) with ReLU activation
    maps the jet representation to four output logits,
    one for each flavour class ($b$, $c$, light, $\tau$).
\end{enumerate}

The model contains approximately \textbf{2.3 million} trainable parameters.
All activations are ReLU; dropout is configurable.
The \texttt{GN2.from\_checkpoint()} class method enables reproducible
reloading of trained models from disk.

% ─── Training ────────────────────────────────────────────────────────────────
\section{Training (\texttt{train.py})}

Training minimises cross-entropy loss on the four-class jet-flavour
classification task.
Jets carrying label $-1$ (unmapped flavour codes) are excluded via
\texttt{ignore\_index=-1}.
The \textbf{AdamW} optimiser is used with weight decay
$\lambda = 10^{-5}$.
The learning-rate schedule exactly reproduces the published ATLAS recipe:
linear warmup from $10^{-7}$ to $5 \times 10^{-4}$ over the first 1\% of
training steps, followed by cosine annealing down to $10^{-5}$.
Gradient norms are clipped to 1.0 at each step.
On CUDA devices, automatic mixed precision (AMP) is activated transparently.
The best checkpoint by validation loss is saved to disk and reloaded at the
end of training; per-epoch losses and learning rates are streamed to
TensorBoard.

% ─── Evaluation ──────────────────────────────────────────────────────────────
\section{Evaluation (\texttt{evaluate.py})}

At evaluation time the best checkpoint is loaded and run on the held-out test
set.
The pipeline produces the following outputs:

\begin{itemize}[leftmargin=*, itemsep=2pt, topsep=2pt]
  \item \texttt{metrics.json} --- overall accuracy and per-class
    precision, recall, and F1-score.
  \item \texttt{confusion\_matrix.pdf} --- normalised confusion matrix.
  \item \texttt{score\_distributions.pdf} --- softmax score distributions
    split by true jet flavour.
  \item \texttt{roc\_db.pdf} --- ROC curves (background rejection vs.\
    $b$-jet efficiency) for the discriminant
    $D_b = \log\!\bigl(p_b / [f_c p_c + f_\tau p_\tau +
    (1-f_c-f_\tau) p_u]\bigr)$,
    with ATLAS defaults $f_c = 0.20$, $f_\tau = 0.05$.
  \item \texttt{roc\_dc.pdf} --- analogous ROC curves for the $c$-tagging
    discriminant $D_c$ ($f_b = 0.30$, $f_\tau = 0.01$).
\end{itemize}

All figures use the \texttt{mplhep} ATLAS style.

% ─── Results ─────────────────────────────────────────────────────────────────
\section{Results}

The table below reports the background rejection values measured by the ATLAS
Collaboration for GN2 at a 70\% $b$-jet tagging efficiency in simulated
$t\bar{t}$ events, as reference benchmarks.
Quantitative results obtained with this implementation will be inserted below
upon completion of the training run.

\vspace{6pt}
\begin{center}
\begin{tabular}{lcc}
  \toprule
  Background & GN2 rejection (ATLAS) & This implementation \\
  \midrule
  $c$-jet    & $> 3\times$ vs.\ DL1d   & \textit{[to be filled]} \\
  light-jet  & $> 1.6\times$ vs.\ DL1d & \textit{[to be filled]} \\
  $\tau$-jet & up to $8\times$ vs.\ DL1d & \textit{[to be filled]} \\
  \bottomrule
\end{tabular}
\end{center}

\end{document}